%%%%%%%%%%%%%%%%
%%% exod 1 %%%
%%%%%%%%%%%%%%%%

\Chapter{1}

\Verse{1}
{
	\hR{וְאֵלֶּה שְׁמוֹת בְּנֵי יִשְׂרָאֵל הַבָּאִים מִצְרָיְמָה אֵת יַעֲקֹב אִישׁ וּבֵיתוֹ בָּאוּ׃}
}
{
	\eR{
		And these are the names of the children of Israel who came
		to Egypt. With Jacob, each and his household had come:
	}
}
{
	Wake up sleepers\footnote{Original Text (OT) has Heb. נרדם here.}
	
	Prop\footnote{Spelling varies between manuscripts, most notably `אפ יתר' in OT.}
	yourselves\footnote{Lit. `יהו selves' in OT.}
	up \& listen.\footnote{Lit. `\[н\]` in OT, believed to be a paronomasia of something like\
						   mixed Cyrillic and programmer, though scholarly consensus is mixed to say the least.}

	He\footnote{Lit. `T\=a' in OT.} who has ears\footnote{Lit. `\chineseB{一二}s'.} let him hear.\footnote{Lit. `He耳'.}

}

\Verse{2}
{
	\hR{רְאוּבֵן שִׁמְעוֹן לֵוִי וִיהוּדָה׃}
}
{
	\eR{Reuben, Simeon, Levi, and Judah,}
}
{
}

\Verse{3}
{
	\hR{יִשָּׂשכָר זְבוּלֻן וּבִנְיָמִן׃}
}
{
	\eR{Issachar, Zebulun, and Benjamin,}
}
{
}

\Verse{4}
{
	\hR{דָּן וְנַפְתָּלִי גָּד וְאָשֵׁר׃}
}
{
	\eR{Dan and Naphtali, Gad and Asher.}
}
{
}

\Verse{5}
{
	\hR{וַיְהִי כּל נֶפֶשׁ יֹצְאֵי יֶרֶךְ יַעֲקֹב שִׁבְעִים נָפֶשׁ וְיוֹסֵף הָיָה בְמִצְרָיִם׃}
}
{
	\eR{
		And all the persons coming out from Jacob’s thigh were
		seventy persons. And Joseph had been in Egypt.
	}
}
{
%	\footnote{
%		seventy. Septuagint and Qumran texts read “seventy-five.” The difference
%		of five reflects Joseph’s sons Ephraim and Manasseh plus Manasseh’s son
%		and Ephraim’s two sons. These are named in Gen 46:20 in the Septuagint
%		but not in the Masoretic Text. The problem is that the children of
%		Israel are identified here in v. 1 as those “who came to Egypt,” which
%		would exclude Joseph’s children and grandchildren because they do not
%		come to Egypt; they are born there. But the people are identified in v.
%		5 as “all the persons coming out from Jacob’s thigh,” which would
%		include Joseph’s offspring no matter where they were born. It appears
%		that each scribe used the number that fit what he understood to be the
%		logic of the text. (Small differences of a few letters in the text such
%		as this occur throughout the Torah. They show that schemes that count
%		the letters of the Bible in order to uncover hidden “codes” are
%		nonsense. Such codes became popular recently but were carried out by
%		persons who were unqualified to do this serious study of the text of the
%		Torah in all its versions. Examination of their misuse of statistical
%		analysis further undermined their claims, which unfortunately misled
%		sincere people. Let us hope that the repudiation of such popular schemes
%		will serve to remind people anew that what is great in the Torah is in
%		what it says, not in secret patterns in its letters.)
%	}
}

\Verse{6}
{
	\hJ{וַיָּמת יוֹסֵף וְכל אֶחָיו וְכֹל הַדּוֹר הַהוּא׃}
}
{
	\eJ{
		And Joseph and all of his brothers and all of that
		generation died.
	}
}
{
}

\Verse{7}
{
	\hP{וּבְנֵי יִשְׂרָאֵל פָּרוּ וַיִּשְׁרְצוּ וַיִּרְבּוּ וַיַּעַצְמוּ בִּמְאֹד מְאֹד וַתִּמָּלֵא הָאָרֶץ אֹתָם׃}
}
{
	\eP{
		And the children of Israel were fruitful and teemed and
		multiplied and became very, very powerful, and the land
		was filled with them.
	}
}
{
}

\Verse{8}
{
	\hE{וַיָּקם מֶלֶךְ חָדָשׁ עַל מִצְרָיִם אֲשֶׁר לֹא יָדַע אֶת יוֹסֵף׃}
}
{
	\eE{And a new king rose over Egypt—who did not know Joseph.}
}
{
%	\footnote{
%		a new king. There are five Pharaohs in the Torah: the Pharaoh who
%		thought Sarah was Abraham’s sister, the Pharaoh who knew Joseph, the
%		Pharaoh who did not know Joseph, the Pharaoh who sought to kill Moses
%		(who may or may not be the same Pharaoh who did not know Joseph), and
%		the Pharaoh of the exodus. Why are none of their names given? Names of
%		Pharaohs (Shishak, Neco) are given in later books. Their absence in the
%		Torah gives the narrative a nonhistorical quality, which is contrary to
%		the manifest aim of the Torah to present history. One might argue that
%		this is evidence that the stories are not true, that they were made up
%		by writers who could not name these kings because they had no idea of
%		the names of ancient Pharaohs. In the case of the two Pharaohs in
%		Genesis, we have hardly any evidence to argue for or against this. But
%		in the case of the Exodus Pharaohs, I think that there is sufficient
%		likelihood that the oppression and exodus are historical, so there must
%		be some other reason why the Pharaohs are not named. My friend Jonathan
%		Saville suggests that perhaps the reason, consciously or not, is to
%		downgrade the Pharaoh, as when people sometimes avoid saying the name of
%		someone toward whom they feel hostile. Or perhaps the names of the
%		Pharaohs were no longer preserved in the tradition by the time the
%		stories came to be written.
%	}
}

\Verse{9}
{
	\hE{וַיֹּאמֶר אֶל עַמּוֹ הִנֵּה עַם בְּנֵי יִשְׂרָאֵל רַב וְעָצוּם מִמֶּנּוּ׃}
}
{
	\eE{
		And he said to his people, “Here, the people of the
		children of Israel is more numerous and powerful than we.
	}
}
{
}

\Verse{10}
{
	\hE{הָבָה נִתְחַכְּמָה לוֹ פֶּן יִרְבֶּה וְהָיָה כִּי תִקְרֶאנָה מִלְחָמָה וְנוֹסַף גַּם הוּא עַל שֹׂנְאֵינוּ וְנִלְחַם בָּנוּ וְעָלָה מִן הָאָרֶץ׃}
}
{
	\eE{
		Come on, let’s be wise toward it or else it will increase;
		and it will be, when war will happen, that it, too, will
		be added to our enemies and will war against us and go up
		from the land.”
	}
}
{
%	\footnote{
%		will be added. The Hebrew () is punning on the name Joseph (), meaning
%		“may He add.” The Pharaoh does not know Joseph, but when he is pictured
%		as worrying that the people “will be added” this summons Joseph back to
%		mind. This pun is not merely wordplay for its own sake. The notation
%		that this is a Pharaoh who does not know Joseph makes a strong break at
%		the beginning of Exodus from what has come before this in Genesis. The
%		reminder of Joseph’s presence reconnects this phase of the story to
%		everything that has come before it. Even stronger connections are
%		coming.
%	}
}

\Verse{11}
{
	\hE{וַיָּשִׂימוּ עָלָיו שָׂרֵי מִסִּים לְמַעַן עַנֹּתוֹ בְּסִבְלֹתָם וַיִּבֶן עָרֵי מִסְכְּנוֹת לְפַרְעֹה אֶת פִּתֹם וְאֶת רַעַמְסֵס׃}
}
{
	\eE{
		And they set commanders of work-companies over it in order
		to degrade it with their burdens. And they built storage
		cities for Pharaoh: Pithom and Rameses.
	}
}
{
%	\footnote{
%		work-companies. The Hebrew term, missîm, refers not to individually
%		owned household slaves but to a policy of forced labor imposed on an
%		entire community (a corvée). The Israelites build whole cities, and they
%		all live in a particular region of Egypt (Goshen), separate from the
%		Egyptian population (Exod 8:18; 9:26). Centuries later, King Solomon
%		imposes missîm on Israel, requiring, in addition to monetary taxes, a
%		period of labor on national projects. This so infuriates the Israelites
%		that they stone to death the king’s minister of missîm (1 Kings 5:27–28;
%		12:18). Israelites will bear taxation, but the requirement of forced
%		labor implies control over people’s bodies by the government. This is
%		appalling to a people whose recollection of having been slaves is a
%		central doctrine to their understanding of themselves and their history.
%	}
}

\Verse{12}
{
	\hE{וְכַאֲשֶׁר יְעַנּוּ אֹתוֹ כֵּן יִרְבֶּה וְכֵן יִפְרֹץ וַיָּקֻצוּ מִפְּנֵי בְּנֵי יִשְׂרָאֵל׃}
}
{
	\eE{
		And the more they degraded it, the more it increased, and
		the more it expanded; and they felt a disgust at the
		children of Israel.
	}
}
{
}

\Verse{13}
{
	\hP{וַיַּעֲבִדוּ מִצְרַיִם אֶת בְּנֵי יִשְׂרָאֵל בְּפָרֶךְ׃}
}
{
	\eP{
		And Egypt made the children of Israel serve with
		harshness;
	}
}
{
}

\Verse{14}
{
	\hP{וַיְמָרְרוּ אֶת חַיֵּיהֶם בַּעֲבֹדָה קָשָׁה בְּחֹמֶר וּבִלְבֵנִים וּבְכל עֲבֹדָה בַּשָּׂדֶה אֵת כּל עֲבֹדָתָם אֲשֶׁר עָבְדוּ בָהֶם בְּפָרֶךְ׃}
}
{
	\eP{
		and they made their lives bitter with hard work, with
		mortar and with bricks and with all work in the field—all
		their work that they did for them—with harshness.
	}
}
{
}

\Verse{15}
{
	\hE{וַיֹּאמֶר מֶלֶךְ מִצְרַיִם לַמְיַלְּדֹת הָעִבְרִיֹּת אֲשֶׁר שֵׁם הָאַחַת שִׁפְרָה וְשֵׁם הַשֵּׁנִית פּוּעָה׃}
}
{
	\eE{
		And the king of Egypt said to the Hebrew midwives—of whom
		the name of one was Shiphrah and the name of the second
		was Puah—
	}
}
{
%	\footnote{
%		Hebrew midwives. The Hebrew may be read as “Hebrew midwives,” meaning
%		that these two women are Israelites; or it may be “midwives of the
%		Hebrews,” in which case one cannot know whether or not they themselves
%		are Israelites. “Hebrew midwives” is more likely because the Israelites
%		are never referred to as “the Hebrews” by the narrator in the Torah. It
%		is a term used in quotation when speaking to foreigners (Gen 40:15; Exod
%		5:3; see the comment on Gen 14:13) or as an adjective in the fixed
%		phrase “Hebrew slave” (see the comment on Exod 21:2). Here, too, it may
%		be such an adjectival usage. This is supported by the fact that these
%		names are much more likely to be Semitic than Egyptian, implying that
%		the midwives are Israelites.
%	}
}

\Verse{16}
{
	\hE{וַיֹּאמֶר בְּיַלֶּדְכֶן אֶת הָעִבְרִיּוֹת וּרְאִיתֶן עַל הָאבְנָיִם אִם בֵּן הוּא וַהֲמִתֶּן אֹתוֹ וְאִם בַּת הִוא וָחָיָה׃}
}
{
	\eE{
		and he said, “When you deliver the Hebrew women, and you
		look at the two stones, if it’s a boy then kill him, and
		if it’s a girl then she’ll live.”
	}
}
{
%	\footnote{
%		the two stones. This is often understood to mean some sort of birthing
%		stool made of two stones, but the more natural understanding here in the
%		context of identifying boys is that the two stones refer to the
%		testicles.
%	}
}

\Verse{17}
{
	\hE{וַתִּירֶאןָ הַמְיַלְּדֹת אֶת הָאֱלֹהִים וְלֹא עָשׂוּ כַּאֲשֶׁר דִּבֶּר אֲלֵיהֶן מֶלֶךְ מִצְרָיִם וַתְּחַיֶּיןָ אֶת הַיְלָדִים׃}
}
{
	\eE{
		And the midwives feared {\God} and did not do what the king
		of Egypt had spoken to them, and they kept the boys alive.
	}
}
{
}

\Verse{18}
{
	\hE{וַיִּקְרָא מֶלֶךְ מִצְרַיִם לַמְיַלְּדֹת וַיֹּאמֶר לָהֶן מַדּוּעַ עֲשִׂיתֶן הַדָּבָר הַזֶּה וַתְּחַיֶּיןָ אֶת הַיְלָדִים׃}
}
{
	\eE{
		And the king of Egypt called the midwives and said to
		them, “Why have you done this thing and kept the boys
		alive?”
	}
}
{
}

\Verse{19}
{
	\hE{וַתֹּאמַרְןָ הַמְיַלְּדֹת אֶל פַּרְעֹה כִּי לֹא כַנָּשִׁים הַמִּצְרִיֹּת הָעִבְרִיֹּת כִּי חָיוֹת הֵנָּה בְּטֶרֶם תָּבוֹא אֲלֵהֶן הַמְיַלֶּדֶת וְיָלָדוּ׃}
}
{
	\eE{
		And the midwives said to Pharaoh, “Because the Hebrews
		aren’t like the Egyptian women, because they’re animals!
		Before the midwife comes to them, they’ve given birth!”
	}
}
{
%	\footnote{
%		they’re animals! The vowels inserted in the Masoretic Text would make
%		this an adjective (“they’re lively”), but that form of the word does not
%		occur anywhere else in the Bible. I think that it is more likely that
%		the midwives are meant to be speaking in this negative way about the
%		Israelite women in order to hide their own violation of the king’s
%		order.
%	}
}

\Verse{20}
{
	\hE{וַיֵּיטֶב אֱלֹהִים לַמְיַלְּדֹת וַיִּרֶב הָעָם וַיַּעַצְמוּ מְאֹד׃}
}
{
	\eE{
		And {\God} was good to the midwives. And the people
		increased, and they became very powerful,
	}
}
{
}

\Verse{21}
{
	\hE{וַיְהִי כִּי יָרְאוּ הַמְיַלְּדֹת אֶת הָאֱלֹהִים וַיַּעַשׂ לָהֶם בָּתִּים׃}
}
{
	\eE{
		and it was because the midwives feared {\God}, and He made
		them households.
	}
}
{
}

\Verse{22}
{
	\hJ{וַיְצַו פַּרְעֹה לְכל עַמּוֹ לֵאמֹר כּל הַבֵּן הַיִּלּוֹד הַיְאֹרָה תַּשְׁלִיכֻהוּ וְכל הַבַּת תְּחַיּוּן׃}
}
{
	\eJ{
		And Pharaoh commanded all of his people, saying, “Every
		son who is born: you shall throw him into the Nile. And
		every daughter you shall keep alive.”
	}
}
{
}
